%
% 32-nullstellenq.tex -- Nullstellen von Polynomen über Q
%
% (c) 2026 Prof Dr Andreas Müller
%
\subsection{Nullstellen von Polynomen über $\mathbb{Q}$}
Der Fundamentalsatz der Algebra löst das Problem der Faktorisierung
ein einem einzigen Körper.
Jedes Polynom $p(Z)\in \mathbb{C}[Z]$ lässt sich in Linearfaktoren
zerlegen.
Dies ermöglicht zum Beispiel die Partialbruchzerlegung rationaler
Funktionen und steht damit am Anfang der Lösung des Integrationsproblem
für rationale Funktionen, welches in Kapitel~\ref{chapter:ratint}
gelöst wird.
\index{Partialbruchzerlegung}%
Jede rationale Funktion mit komplexen Koeffizienten hat eine
Stammfunktion, die sich als rationale Funktion und eine Linearkombination
von Logarithmusfunktionen der Form $\log(z-\alpha_i)$ ist, wobei die
$\alpha_i\in\mathbb{C}$ die Nullstellen des Nennerpolynoms sind.

Trotz dem unbestreitbaren Erfolg der komplexen Zahlen bleibt ein
grundlegendes Problem für die Anwendung in der Computeralgebra.
\index{Computeralgebra}%
Nicht jede komplexe Zahl hat eine exakte Darstellung in einem Computersystem.
Da $\mathbb{C}=\mathbb{R}(i)$ eine algebraische Erweiterung der 
überabzählbar unendlichen Menge der reellen Zahlen ist, gibt es
überabzählbar unendlich viele komplexe Zahlen.
In einem Computer kann man aber prinzipiell nur abzählbar unendlich
viele Zahlen exakt darstellen.
Gleitkomma-Arithmetik löst das Problem dadurch, dass reelle Zahlen
\index{Gleitkomma-Arithmetik}%
und damit auch komplexe Zahlen approximiert werden.
Die meisten Zahlen, die in numerischen Rechnungen auftreten, sind
daher nur Approximationen und man kann nicht erwarten, die Nullstellen
eines Polynoms exakt zu bestimmen, selbst bei Polynomen mit
rationalen Koeffizienten, die eine endliche Spezifikation haben.

Ein Computeralgebrasystem möchte entscheiden können, ob eine Gleichung
exakt erfüllt ist.
Dazu ist eine Approximation nicht ausreichen.
Es muss ein Zahlenbereich gefunden werden, in dem solche exakten Aussagen
möglich sind.
Er muss die arithmetische Grundoperationen ermöglichen, d.~h.~er muss
die rationalen Zahlen enthalten.
Die verbreitete freie Softwarebibliothek GMP \cite{buch:gmp} realisiert
\index{GMP}%
die Arithmetik mit rationalen Zahlen mit beliebig grossen Zählern und
Nennern, solange diese im Speicher Platz haben.
Man kann zwar nicht jede rationale Zahl im Computer darstellen, aber
man kann jede Zahl, die man im Laufe einer Berechnung antrifft, exakt
darstellen.

In einem Computeralgebrasystem muss man also nicht für jedes mathematische
Objekt von Anfang an eine exakte Darstellung haben, aber man muss
jedes Objekt, das im Laufe einer Rechnunung auftritt, exakt darstellen
können.
Möchte man eine quadratische Gleichung lösen, braucht es im Allgemeinen
eine Darstellung für jede beliebige Quadratwurzel von rationalen
Zahlen.
Ein Computeralgebrasystem führt dazu neue Symbole mit geeigneten
Rechenregeln ein.
Um mit der Wurzel $\!\sqrt{2}$ rechnen zu können, muss man dem
Computeralgebrasystem ein neues Objekt \texttt{sqrt(2)} hinzufügen,
welches die Eigenschaft $\texttt{sqrt(2)}^2 = 2$ hat.
Diese Konstruktion kann iteriert werden, \text{sqrt(sqrt(2))} ist
ein Objekt welches die algebraischen Eigenschaften einer vierten Wurzel
hat.

Der Fundamentalsatz der Algebra erhält in diesem Zusammenhang eine
neue Rolle.
Er garantiert, dass jedes komplexe Polynome eine Nullstelle hat, die
man bei Bedarf auch numerisch berechnen kann.
Für eine algebraisch exakte Darstellung ist diese numerische Darstellung
aber nicht nötig.
Viel wichtiger es, eine solche Nullstelle als ein Element zu interpretieren,
um das der Körper $\mathbb{Q}$ erweitert werden muss, damit die Rechnung
weitergeführt werden kann.

Die Berechnung in einem Computeralgebrasystem beginnt also mit den rationalen
Zahlen $\mathbb{Q}$.
Jedesmal, wenn eine Wurzel oder eine Nullstelle eines Polynoms auftritt,
wird der Körper der darstellbaren Zahlen um das algebraische Element
$\alpha$ erweitert.
Die weiteren Rechnungen finden im erweiterten Körper $K=\mathbb{Q}(\alpha)$
statt.
Die algebraischen Eigenschaften des Elements $\alpha$ sind Teil der
Spezifikation von $\alpha$.
Auch die imaginäre Einheit ist von dieser Art.

Ein Computeralgebrasystem muss aber auch mit anderen Grössen umgehen
können, zum Beispiel mit der Kreiszahl $\pi$ oder der Basis $e$
\index{Kreiszahl}%
\index{pi@$\pi$}%
\index{e@$e$}%
der natürlichen Logarithmen.
Auch diese Zahlen existieren als Elemente in $\mathbb{C}$.
Für die Rechnung in einem Computeralgebrasystem ist die Darstellung von $\pi$
oder $e$ als Gleitkommazahl nicht ausreichend.
Aber auch diese Zahlen können durch ihre algebraischen Eigenschaften
charakterisiert werden.
Es gibt zwar kein Polynom in $\mathbb{Q}[z]$, das $\pi$ oder $e$ als
Nullstelle hat, da diese Zahlen transzendent sind.
Die Potenzen $\pi^k$, $k\in\mathbb{N}$ sind über $\mathbb{Q}$ linear
unabhängig.
Das bedeutet aber nur, dass man eine Potenz von $\pi$ niemals
durch niedrigere Potenzen wird ausdrücken können, wie das bei den
algebraischen Elementen der Fall ist.
Dafür gibt es Beziehungen wie $e^{i\pi}=-1$, mit denen Exponentialausdrücke
vereinfacht werden können, die $\pi$ im Exponenten enthalten.
Für die praktischen Zwecke eines Computeralgebrasystems genügt dies.
Für jeden im Laufe der Rechnung auftretenden Ausdruck lässt sich eine
Darstellung im Computer finden.

In der Analysis lernt man, dass für die Definition der Ableitung einer
Funktion ein Grenzwert in $\mathbb{R}$ nötig ist.
In Abschnitt~\ref{buch:holomorph:section:holomorph} werden wir zeigen,
wie sich die Ableitung auf komplexe Funktionen ausdehnen lässt.
Auch in diesem Szenario garantieren die analytischen Eigenschaften des
Körpers $\mathbb{C}$, dass die Ableitung als Zahl in $\mathbb{C}$ existiert.
Es ist zwar möglich, Ableitungen numerisch zu berechnen, eine besonders
elegante Methode wird in Kapitel~\ref{chapter:step} beschrieben.
Dies ist aber wieder nicht die Art und Weise, wie ein Computeralgebrasystem
mit der Ableiltung umgeht.
Hierfür werden die algebraischen Eigenschaften des Ableitungsoperators
verwendet, die man als die Ableitungsregeln aus der Analysis kennt.
Sie ermöglichen, die Ableitung beliebig komplizierter Ausdrücke
zu berechnen, sofern man nur die Ableitungsregeln für die Komponenten
des Ausdrucks kennt.
Man muss zum Beispiel nicht wissen, wie man einen Wert von $f(x)=e^{1291x}$
berechnen kann, um die Ableitung zu berechnen.
Die Ableitungsregeln sagen $f'(x)=1291 e^{1291x}$.
Abschnitt~\ref{buch:intalg:section:elementar} wird dies formalisieren.

Auch die Definition des Integrals verwendet die analytische
Vollständigkeit von $\mathbb{C}$.
Die Existenz von Werten eines Integrals ist dadurch gesichert.
Ein Computeralgebrasystem ist jedoch nicht an Werten interessiert,
sondern an algebraischen Ausdrücken für die Stammfunktion.
Um einen solchen zu finden, ist die Definition des Riemann-Integrals
\index{Riemann-Integral}%
\index{Integral, Riemann-}%
als Grenzwert einer Summe von Rechteckflächen nicht zweckmässig.
Die Lösung wird durch den Hauptsatz der Infinitesimalrechnung gegeben.
\index{Hauptsatz}%
Die Stammfunktion ist durch die rein algebraische Eigenschaft
charakterisiert, dass der Ableitungsoperator aus der Stammfunktion
den Integranden produziert.
Auch das Integrationsproblem ist damit auf ein rein algebraisches
Problem reduziert, welches in Kapitel~\ref{chapter:intalg}
studiert wird.













